\documentclass{article}
\usepackage{tcolorbox}
\tcbuselibrary{skins} 
\usepackage[a4paper]{geometry}
\usepackage{fancyhdr}
\pagestyle{fancy}
\lhead{Mathe Präsentation \textit{Hilfen}}
\rhead{zum 25. Februar 2026}
 
\newcommand{\task}[2]{% 
\begin{center} 
 \begin{tcolorbox}[width=0.95\linewidth,
  enhanced,
  title=\textbf{#1)}, 
  attach boxed title to top left={
   xshift=0.75cm,
   yshift=-\tcboxedtitleheight/2,
   yshifttext=-\tcboxedtitleheight/2,
  }, 
  boxed title style={
    sharp corners,
    colback=white,
    colframe=black
  },
  coltitle=black, 
  colframe=black, colback=white,
  drop shadow=black!50!white,
  sharp corners, 
  boxrule=1pt, boxsep=8pt,
  left=0pt,right=0pt,top=0pt,bottom=0pt]
  #2
 \end{tcolorbox}
\end{center}  
} 
 
\begin{document}
\noindent \texttt{Jahr 2025, eA, Aufgaben 1B und 2A} 
\section*{Aufgabe 1B}
\task{a}{Wo wären die Enden des Fensters? Was ist die Höhe dort?}
\task{b}{Was ist ein Längenverhältnis? Was ist dessen Formel?}
\task{c}{Wo ist die höchste Steigung? Was ist die Steigung (als Winkel) dort?}
\task{d}{Wie hoch ist der Bereich dann? In welchem Intervall existiert dieser?}
\task{e}{Was ist eine Tangente, welcher Form folgt sie? Welche Bedingungen hat sie hier? Was muss ihre Steigung sein?}
\task{f}{Wofür stehen die einzelnen Variablen und Funktionen?} 
\newpage 
\section*{Aufgabe 2A} 
\task{a}{Was ist bereits gegeben? Was folgt daraus?}
\task{b}{Was sind die jeweiligen einzelnen Wahrscheinlichkeiten? Was heißt \textit{entweder oder}?}
\task{c}{Wie wird so ein Intervall berechnet? Was sind die Variablen hier?}
\task{d}{
\subsubsection*{Lösungsweg I}
Was ist die normale Form der Normalverteilung? Was ist hier demnacht $\mu$ und $\sigma$? Wie werden die Intervallswahrscheinlichkeiten berechnet? 
\subsubsection*{Lösungsweg II}
Wie sind Intervallswahrscheinlichkeiten eigentlich definiert? 
}
\task{e}{Was muss erfüllt sein, damit Bedingung II erfüllt ist? Ist dann Bedingung III auch automatisch erfüllt?} 
\end{document}